% Reconstructed from upstream/Algorithms-JeffE.pdf, PDF pages 371-396 (printed 353-378).
\epigraph{For a long time it puzzled me how something so expensive, so leading edge, could be so
useless, and then it occurred to me that a computer is a stupid machine with the ability
to do incredibly smart things, while computer programmers are smart people with the
ability to do incredibly stupid things. They are, in short, a perfect match.
}{Bill Bryson, \textit{Notes from a Big Country} (1999)}

\epigraph{Shortly after the “iron curtain” fell in 1990, an American and a Russian, who had both
worked on the development of weapons, met. The American asked: “When you
developed the Bomb, how were you able to perform such an enormous amount of
computing with your weak computers?” The Russian responded: “We used better
algorithms.”
}{Yefim Dinitz, in “Dinitz’ Algorithm: The Original Version\\and Even’s Version” (2006)}

\epigraph{So I’m the bad guy now I hear, because I don’t go with the flow.
Don’t ever go with the flow, be the flow.
}{Jay-Z, "Stream of Consciousness", May 16, 2015}

\chapter{Applications of Flows and Cuts}
\label{ch:11}

\section{Edge-Disjoint Paths}
\label{sec:11.1}

One of the easiest applications of maximum flows is computing the maximum number of edge-disjoint paths between two specified vertices $s$ and $t$ in a directed graph $G$ using maximum flows.
A set of paths in $G$ is \textit{edge-disjoint} if each edge in $G$ appears in at most one of the paths; several edge-disjoint paths may pass through the same vertex, however.

If we give each edge capacity $1$, then the maximum flow from $s$ to $t$ pushes either $0$ or $1$ units of flow along each edge.
The flow-decomposition theorem implies that the subgraph $S$ of saturated edges is the union of several edge-disjoint paths and cycles.
Moreover, the number of paths in this decomposition is exactly equal to the value of the flow.
Extracting the actual paths from $S$ is straightforward—follow any directed path in $S$ from $s$ to $t$, remove that path from $S$, and recurse.

Conversely, we can transform any collection of $k$ edge-disjoint paths into a flow by pushing one unit of flow along each path from $s$ to $t$; the value of the resulting flow is exactly $k$.
It follows that any maximum flow algorithm actually computes the largest possible set of edge-disjoint paths.

If we use Orlin’s algorithm to compute maximum flows, we can compute edge-disjoint paths in $O(VE)$ time, but Orlin’s algorithm is overkill for this simple application.
The cut $(\{s\}, V \setminus \{s\})$ has capacity at most $V - 1$, so the maximum flow has value at most $V - 1$.
Thus, Ford and Fulkerson’s original augmenting path algorithm already runs in $O(|f^*|E) = \symbfit{O(VE)}$ \textbf{\textit{time}}.

The same algorithm can also be used to find edge-disjoint paths in \textit{undirected} graphs.
First, replace every undirected edge $uv$ in $G$ with a pair of directed edges $\arc{u}{v}$ and $\arc{v}{u}$, each with unit capacity, and call the resulting directed graph $G'$.
Next, compute a maximum $(s,t)$-flow $f^*$ in $G'$ using Ford-Fulkerson.
For any edge $uv$ in $G$, if $f^*$ saturates both directed edges $\arc{u}{v}$ and $\arc{v}{u}$ in $G'$, we can remove \textit{both} edges from the flow without changing its value.
(More generally, we can find an \textit{acyclic} maximum flow in $G'$ by canceling all cycles in $f^*$, not only cycles of length $2$.)
Thus, without loss of generality, $f^*$ assigns a unique direction to each saturated edge.
Finally, we can extract the edge-disjoint paths by searching the subgraph of directed edges saturated by $f^*$.

\section{Vertex Capacities and Vertex-Disjoint Paths}
\label{sec:11.2}

Now suppose the vertices of the input graph $G$ have capacities, not just the edges.
In addition to our other constraints, for each vertex $v$ other than $s$ and $t$, we require the total flow into $v$ (and therefore the total flow out of $v$) to be at most some non-negative value $c(v)$:
\[
\sum_{\arc{u}{v}} f(\arc{u}{v}) \le c(v).
\]
Can we still compute maximum flows with these new vertex constraints?

In 1962, Ford and Fulkerson proposed the following reduction to a flow network $\bar{G}$ with only edge capacities.
Replace every vertex $v$ with two vertices $v_{\textup{in}}$ and $v_{\textup{out}}$, connected by an edge $\arc{v_{\textup{in}}}{v_{\textup{out}}}$ with capacity $c(v)$, and then replace every directed edge $\arc{u}{v}$ with the edge $\arc{u_{\textup{out}}}{v_{\textup{in}}}$ (keeping the same capacity).
Routine definition-chasing implies that every feasible $(s_{\textup{out}}, t_{\textup{in}})$-flow in $\bar{G}$ is equivalent to a feasible $(s,t)$-flow with the same value in the original graph $G$, and vice versa.
In particular, every maximum flow in $\bar{G}$ is equivalent to a maximum flow in $G$.
The reduction from $G$ to $\bar{G}$ takes $O(E)$ time, after which we can compute the maximum flow in $\bar{G}$ using Orlin’s algorithm.
Altogether, computing the maximum flow in $G$ requires $\symbfit{O(VE)}$ \textbf{\textit{time}}.

It is now easy to compute the maximum number of \textit{vertex}-disjoint paths from $s$ to $t$ in $O(VE)$ time: Assign capacity $1$ to every vertex and compute a maximum flow!

\begin{figure}[ht]
\centering
\includegraphics{p373-Im157}
\caption{Reducing vertex-disjoint paths in $G$ to edge-disjoint paths in $\bar{G}$.}
\label{fig:11.1}
\end{figure}

\section{Bipartite Matching}
\label{sec:11.3}

Another natural application of maximum flows is finding maximum \textbf{\textit{matchings}} in bipartite graphs.
A matching is a subgraph in which every vertex has degree at most one, or equivalently, a collection of edges such that no two share a vertex.
The problem is to find a matching with the maximum number of edges.

For example, suppose we have a set of doctors who are looking for jobs, and a set of hospitals who are looking for doctors.
Each doctor lists all hospitals where they are willing to work, and each hospital lists all doctors they are willing to hire.
Our task is to find the largest subset of doctor-hospital hires that everyone is willing to accept.\footnote{This problem is very different from the stable matching problem we saw in \chapref{ch:4}, because we aren’t trying to make each doctor and hospital as happy as possible.}
This problem is equivalent to finding a maximum matching in a bipartite graph whose vertices are the doctors and hospitals, and there is an edge between a doctor and a hospital if and only if each find the other acceptable.

We can solve this problem by reducing it to a maximum flow problem, as follows.
Let $G$ be the given bipartite graph with vertex set $L \cup R$, such that every edge joins a vertex in $L$ to a vertex in $R$.
We create a new \textit{directed} graph $G'$ by (1) orienting each edge from $L$ to $R$, (2) adding a new source vertex $s$ with edges to every vertex in $L$, and (3) adding a new target vertex $t$ with edges from every vertex in $R$.
Finally, we assign every edge in $G'$ a capacity of $1$.

Any matching $M$ in $G$ can be transformed into a flow $f_M$ in $G'$ as follows: For each edge $uw$ in $M$, push one unit of flow along the path $\arc{\arc{\arc{s}{u}}{w}}{t}$.
These paths are disjoint except at $s$ and $t$, so the resulting flow satisfies the capacity constraints.
Moreover, the value of the resulting flow is equal to the number of edges in $M$.

Conversely, consider any $(s,t)$-flow $f$ in $G'$, computed using the Ford-Fulkerson augmenting path algorithm.
Because the edge capacities are integers, the Ford-Fulkerson algorithm assigns an integer flow to every edge.
(This is easy to verify by induction, hint, hint.)
Moreover, since each edge has \textit{unit} capacity, the computed flow either saturates ($f(e) = 1$) or avoids ($f(e) = 0$) every edge in $G'$.
Finally, since at most one unit of flow can enter any vertex in $U$ or leave any vertex in $W$, the saturated edges from $U$ to $W$ form a matching in $G$.
The size of this matching is exactly $|f|$.

Thus, the size of the maximum matching in $G$ is equal to the value of the maximum flow in $G'$, and provided we compute the maxflow using augmenting paths, we can convert the actual maxflow into a maximum matching in $O(E)$ time.
We can compute the maximum flow in $\symbfit{O(VE)}$ \textbf{\textit{time}} using either Orlin’s algorithm or off-the-shelf Ford-Fulkerson.

\begin{figure}[ht]
\centering
\includegraphics{p374-Im158}\hfil\includegraphics{p374-Im159}
\caption{A maximum matching in a bipartite graph $G$, and the corresponding maximum flow in $G'$.}
\label{fig:11.2}
\end{figure}

It is enlightening to interpret the behavior of Ford-Fulkerson in $G'$ in terms of the original bipartite graph $G$.
The algorithm maintains a matching $M$ in $G$, which is initially empty; the edges of $M$ correspond to edges in $G'$ that carry flow.
Call a vertex of $G$ \textit{matched} if is an endpoint of some edge in $M$, and \textit{unmatched} otherwise.
In each iteration of the algorithm, we look for an \textit{alternating path} in $G$—a path from an unmatched vertex of $L$ to an unmatched vertex in $R$ that alternates between edges in $M$ and edges not in $M$.
(Alternating paths in $G$ correspond exactly to augmenting paths for the current flow in $G'$.)
If we find an augmenting path $P$, we update $M$ to the symmetric difference $M \oplus P$, which increases the number of edges in $M$ by $1$, and continue to the next iteration.
If there is no alternating path, the maxflow-mincut theorem implies that $M$ is a maximum matching, so the algorithm ends.
Finding a single alternating path requires $O(E)$ time, and the algorithm halts after at most $V$ iterations, so the overall algorithm runs in $\symbfit{O(VE)}$ \textbf{\textit{time}}.
\figref{fig:11.3} shows this algorithm in action.

This characterization of maximum bipartite matchings in terms of alternating paths was proved by Claude Berge in 1957 (independently of the maxflow-mincut theorem), although it was already implicit in algorithms described by Harald Kuhn in 1955, by Dénes Kőnig in 1916, and by Carl Jacobi around 1836.

A more sophisticated algorithm, proposed by John Hopcroft and Richard Karp in 1973, computes maximum matchings in bipartite graphs in only $O(\sqrt{V}E)$ time, by finding several disjoint alternating paths in each iteration.

\begin{figure}[htp]
\centering
\includegraphics{p375-Im160}
\caption{An increasing sequence of matchings connected by alternating paths.}
\label{fig:11.3}
\end{figure}

\section{Tuple Selection}
\label{sec:11.4}

The bipartite maximum matching problem is the simplest example of a broader class of problems that I call \textbf{\textit{tuple}} \textbf{\textit{selection}}.\footnote{I couldn’t find a standard name for these problems, so I made up my own. These are sometimes called “assignment problems”, but it’s more common for the phrase “the assignment problem” to refer to the problem of finding a maximum-\textit{weight} bipartite matching in an edge-weighted bipartite graph.}
The input to a tuple selection problem consists of several finite sets $X_1, X_2, \dots, X_d$, each representing a different discrete resource.
Our task is to select the largest possible set of $d$-tuples, each containing exactly one element from each set $X_i$, subject to several \textit{capacity} constraints of the following form:
\begin{itemize}
\item For each index $i$, Each element $x \in X_i$ can appear in at most $c(x)$ selected tuples.
\item For each index $i$, any two elements $x \in X_i$ and $y \in X_{i+1}$ can appear in at most $c(x,y)$ selected tuples.
\end{itemize}
\noindent Each of the upper bounds $c(x)$ and $c(x,y)$ is either a (typically small) non-negative integer or $\infty$.

In the maximum-matching problem, we have $d = 2$ resources, each element $x$ has capacity $c(x) = 1$, and each pair $(x,y)$ has capacity $c(x,y) = 1$ or $c(x,y) = 0$, depending on whether or not $xy$ is an edge in the underlying bipartite graph.

Because the resources are linearly ordered, and only pairs of objects in \textit{adjacent} subsets $X_i$ and $X_{i+1}$ are constrained,\footnote{If pairs of objects from even one non-adjacent pair of subsets ($X_i$ and $X_j$ where $j > i + 1$) are also constrained, the problem becomes NP-hard, by a straightforward reduction from \Proc{Exact-3DimensionalMatching}. We’ll discuss NP-hardness in the next chapter.} the tuple selection problem can be reduced to a maximum-flow problem in a directed graph $G$ defined as follows:
\begin{itemize}
\item $G$ contains a vertex for each element of each set $X_i$, as well as a source vertex $s$ and a target vertex $t$. Each vertex $x$ (except $s$ and $t$) has capacity $c(x)$.
\item $G$ contains an edge $\arc{s}{w}$ for each element $w \in X_1$, an edge $\arc{z}{t}$ for each element $z \in X_d$, and an edge $\arc{x}{y}$ with capacity $c(x,y)$ for each pair of elements $x \in X_i$ and $y \in X_{i+1}$, for all $i$. (Optionally, we can omit edges $\arc{x}{y}$ with $c(x,y) = 0$.)
\end{itemize}
\noindent Every path from $s$ to $t$ in $G$ corresponds to (or “is”) a $d$-tuple that we \textit{could} select; conversely, every selectable $d$-tuple that satisfies the stated constraints corresponds to (or “is”) a path from $s$ to $t$ in $G$.

\begin{figure}[ht]
\centering
\includegraphics{p376-Im161}
\caption{The flow network for a tuple selection problem.}
\label{fig:11.4}
\end{figure}

More generally, let $f$ be an arbitrary feasible integer $(s,t)$-flow in $G$.
Because all capacities are integers or $\infty$, the Flow Decomposition Theorem implies that $f$ is the sum of $|f|$ paths from $s$ to $t$, each carrying exactly one unit of flow.
Straightforward definition-chasing implies that the resulting set of tuples satisfies all the capacity constraints.
Conversely, for any set of $k$ tuples that satisfies the capacity constraints, the sum of the $k$ corresponding paths is a feasible integer $(s,t)$-flow with value $k$.

Thus, we can select the maximum number of tuples that satisfy the given capacity constraints by computing a maximum $(s,t)$-flow $f^*$ in $G$ and then computing a flow decomposition of $f^*$.
Because all finite capacities in $G$ are integers, we can assume without loss of generality that $f^*$ is an integer flow, and therefore (by the previous paragraph) corresponds to a valid set of $|f^*|$ tuples.

\subsection{Exam Scheduling}

The following “real world” scheduling problem might help clarify our general reduction.

Sham-Poobanana University has hired you to write an algorithm to schedule their final exams.
There are $n$ different classes, each of which needs to schedule a final exam in one of $r$ rooms during one of $t$ different time slots.
At most one class’s final exam can be scheduled in each room during each time slot; conversely, classes cannot be split into multiple rooms or multiple times.
Moreover, each exam must be overseen by one of $p$ proctors.\footnote{or as they are better known outside the US, invigilators}
Each proctor can oversee at most one exam at a time; each proctor is available for only certain time slots; and no proctor is allowed oversee more than $5$ exams total.
The input to the scheduling problem consists of three arrays:
\begin{itemize}
\item An integer array $E[1\,..\,n]$ where $E[i]$ is the number of students enrolled in the $i$th class.
\item An integer array $S[1\,..\,r]$, where $S[j]$ is the number of seats in the $j$th room. The $i$th class’s final exam \textit{can} be held in the $j$th room if and only if $E[i] \le S[j]$.
\item A boolean array $A[1\,..\,t,1\,..\,p]$ where $A[k,\ell] = \Const{True}$ if and only if the $\ell$th proctor is available during the $k$th time slot.\footnote{Arguably, this information is better represented as a graph, but I thought that would make the reduction more confusing.}
\end{itemize}
\noindent let $N = n + r + tp$ denote the total size of the input.
Your job is to design an algorithm that either schedules a room, a time slot, and a proctor for every class’s final exam, or correctly reports that no such schedule is possible.

This is a standard tuple-selection problem with four resources: classes, rooms, time slots, and proctors.
To solve this problem, we construct a flow network $G$ with six types of vertices—a source vertex $s'$, a vertex $c_i$ for each class, a vertex $r_j$ for each room, a vertex $t_k$ for each time slot, a vertex $p_\ell$ for each proctor, and a target vertex $t'$—and five types of edges, as shown in \figref{fig:11.5}:

\begin{figure}[ht]
\centering
\includegraphics{p377-Im162}
\caption{A flow network for the exam scheduling problem.}
\label{fig:11.5}
\end{figure}

\begin{itemize}
\item An edge $\arc{s'}{c_i}$ with capacity $1$ for each class $i$. (“Each class can hold at most one final exam.”)
\item An edge $\arc{c_i}{r_j}$ with capacity $\infty$ for each class $i$ and room $j$ such that $E[i] \le S[j]$. (“Class $i$ can hold exams in room $j$ if and only if the room has enough seats.”) This is the only place where the enrollments $E[i]$ and seat numbers $S[j]$ are used.
\item An edge $\arc{r_j}{t_k}$ with capacity $1$ for each room $j$ and time slot $k$. (“At most one exam can be held in room $j$ at time $k$.”)
\item An edge $\arc{t_k}{p_\ell}$ with capacity $1$ for time slot $k$ and proctor $\ell$ such that $A[\ell,k] = \Const{True}$. (“A proctor can oversee at most one exam at any time, and only during times that they are available.”)
\item An edge $\arc{p_\ell}{t'}$ with capacity $5$ for each proctor $\ell$. (“Each proctor can oversee at most $5$ exams.”)
\end{itemize}
\noindent (I’m calling the source and target vertices $s'$ and $t'$ instead of $s$ and $t$ only because the problem statement already uses the variable $t$ to denote the number of time slots.)
Altogether, $G$ has $n + r + t + p + 2 = O(N)$ vertices and $O(nr + rt + tp) = O(N^2)$ edges.

Each path from $s'$ to $t'$ in $G$ represents a unique valid choice of class, room, time, and proctor for one final exam; specifically, the class fits into the room, and the proctor is available at that time.
Conversely, for each valid choice (class, room, time, proctor), there is a corresponding path from $s'$ to $t'$ in $G$.
Thus, we can construct a valid schedule for the maximum possible number of exams by computing an maximum $(s',t')$-flow $f^*$ in $G$, decomposing $f^*$ into paths from $s'$ to $t'$, and then transcribing each path into a class-room-time-proctor assignment.
If $|f^*| = n$, we can return the resulting schedule; otherwise, we can correctly report that scheduling all $n$ final exams is impossible.

Constructing $G$ from the given input data by brute force takes $O(E)$ time.
We can compute the maximum flow in $O(VE)$ time using either Ford-Fulkerson (because $|f^*| \le n < V$) or Orlin’s algorithm, and we can compute the flow decomposition in $O(VE)$ time.
Thus, the overall algorithm runs in $O(VE) = \symbfit{O(N^3)}$ \textbf{\textit{time}}.

\section{Disjoint-Path Covers}
\label{sec:11.5}

A \textbf{\textit{path cover}} of a directed graph $G$ is a collection of directed paths in $G$ such that every vertex of $G$ lies on \textit{at least} one path.
A \textbf{\textit{disjoint-}}path cover of $G$ is a path cover such that every vertex of $G$ lies on \textit{exactly} one path.
Every directed graph has a trivial disjoint-path cover consisting of several paths of length zero, but that’s boring.
Instead, let’s look for disjoint-path covers that contain as few paths as possible.
This problem is NP-hard in general—a graph has a disjoint-path cover of size $1$ if and only if it contains a Hamiltonian path—but there is an efficient flow-based algorithm for directed \textit{acyclic} graphs.

To solve this problem for a given directed acyclic graph $G = (V,E)$, we construct a new bipartite graph $G' = (V',E')$ as follows.
\begin{itemize}
\item $H$ contains two vertices $v^\flat$ and $v^\sharp$ for every vertex $v$ of $G$.
\item $H$ contains an undirected edge $u^\flat v^\sharp$ for every directed edge $\arc{u}{v}$ in $G$.
\end{itemize}
\noindent (If $G$ is represented as an adjacency matrix, then $G'$ is the bipartite graph represented by the same adjacency matrix!)

\begin{figure}[ht]
\centering
\includegraphics{p379-Im163}
\caption{Reducing minimum disjoint-path cover of a dag to maximum bipartite matching; squares are flat$^\flat$ and diamonds are sharp$^\sharp$.}
\label{fig:11.6}
\end{figure}

Now I claim that $G$ can be covered by $k$ disjoint paths if and only if the new graph $G'$ has a matching of size $V - k$.
As usual, we prove the equivalence in two stages:
\begin{itemize}
\item[$\Leftarrow$] Suppose $G$ has a disjoint path cover $P$ with $k$ paths; think of $P$ as a subgraph of $G$.
Every vertex in $P$ has in-degree either $0$ or $1$; moreover, there is exactly one vertex with in-degree $0$ in each path in $P$.
It follows that $P$ has exactly $V - k$ edges.
Now define a subset $M$ of the edges of $G'$ as follows:
\[
M := \left\{ u^\flat v^\sharp \in E' \;\middle|\; \arc{u}{v} \in P \right\}.
\]
By definition of disjoint-path cover, every vertex of $G$ has at most one incoming edge in $P$ and at most one outgoing edge in $P$.
We conclude that every vertex of $G'$ is incident to at most one edge in $M$; that is, $M$ is a matching of size $V - k$.
\item[$\Rightarrow$] Suppose $G'$ has a matching $M$ of size $V - k$.
We project $M'$ back to $G$ by defining a subgraph $P = (V, M')$, where
\[
M' := \left\{ \arc{u}{v} \in E \;\middle|\; u^\flat v^\sharp \in M \right\}.
\]
By definition of matching, every vertex of $G$ has at most one incoming edge in $P$ and at most one outgoing edge in $P$.
It follows that $P$ is a collection of disjoint directed paths in $G$; since $P$ includes every vertex, $P$ defines an disjoint path cover with $V - k$ edges.
The number of paths in $P$ is equal to the number of vertices in $G$ that have no incoming edge in $M'$.
We conclude that $P$ contains exactly $k$ paths.
\end{itemize}
\noindent It follows immediately that we can find a minimum disjoint-path cover in $G$ by computing a maximum matching in $G'$, using Ford-Fulkerson’s maximum-flow algorithm, in $O(V'E') = \symbfit{O(VE)}$ \textbf{\textit{time}}.

Despite its formulation in terms of dags and paths, this is really a maximum matching problem: We want to \textit{match} as many vertices as possible to distinct \textit{successors} in the graph.
The number of paths required to cover the dag is equal to the number of vertices with no successor.
(And of course, every bipartite maximum matching problem is really a flow problem.)

\subsection{Minimal Faculty Hiring}

Let’s go back to Sham-Poobanana University for another “real-world” scheduling problem.\footnote{For a somewhat more realistic (and less depressing) formulation of this problem, consider airplanes and flights, or buses and bus routes, instead of professors and classes.}
SPU offers several thousand courses every day.
Due to extreme budget cuts, the university needs to significantly reduce the size of its faculty.
However, because students pay tuition (and the university cannot afford lawyers), the university must retain enough professors to guarantee that every class advertised in the course catalog is actually taught.
How few professors can SPU get away with?
Each remaining faculty member will be assigned a sequence of classes to teach on any given day.
The classes assigned to each professor must not overlap; moreover, there must be enough slack in each professor’s schedule for them to walk from one class to the next.
For purposes of this problem, let’s assume that every professor is capable of teaching every class, and that professors will not have office hours, lunches, or bathroom breaks.\footnote{They will, however, be expected to answer student emails as they walk between classes.}

Concretely, suppose there are $n$ classes offered in $m$ different locations.
The input to our problem consists of the following data:
\begin{itemize}
\item An array $C[1\,..\,n]$ of classes, where each class $C[i]$ has three fields: the starting time $C[i].\Var{start}$, the ending time $C[i].\Var{end}$, and the location $C[i].\Var{loc}$.
\item A two-dimensional array $T[1\,..\,m,1\,..\,m]$, where $T[u,v]$ is the time required to walk from location $u$ to location $v$.
\end{itemize}
\noindent We want to find the minimum number of professors that can collectively teach every class, such that whenever a professor is assigned to teach two classes $i$ and $j$ where $C[j].\Var{start} \ge C[i].\Var{start}$, we actually have
\[
C[j].\Var{start} \;\ge\; C[i].\Var{end} + T\bigl[C[i].\Var{loc},\, C[j].\Var{loc}\bigr].
\]

We can solve this problem by reducing it to a disjoint-path cover problem as follows.
We construct a dag $G = (V,E)$ whose vertices are classes and whose edges represent pairs of classes that are scheduled far enough apart to be taught by the same professor.
Specifically, a directed edge $\arc{i}{j}$ indicates that the same professor can teach class $i$ and then class $j$.
It is easy to construct this dag in $O(n^2)$ time by brute force.
Then we find a disjoint-path cover of $G$ using the matching algorithm described above; each directed path in $G$ represents a legal class schedule for one professor.
The entire algorithm runs in $O(n^2 + VE) = \symbfit{O(n^3)}$ \textbf{\textit{time}}.\footnote{If we assume that every time interval $T[u,v]$ is equal,\footnotemark{} this scheduling problem can actually be solved in $O(n \log n)$ time using a simple greedy algorithm.}\footnotetext{Many American universities schedule ten-minute breaks between classes, under the remarkable belief that a human being can walk from any classroom to any other classroom on the same campus in ten minutes. I invite anyone who thinks this belief is realistic to visit my campus and walk from one Siebel Center to the other.}

Despite its initial description in terms of intervals and distances, this is really a maximum matching problem (which means it’s \textit{really} really a maximum-flow problem).
Specifically, we want to \textit{match} as many classes as possible to the \textit{next} class taught by the same professor.
The number of professors we need is equal to the number of classes with no assigned successor; each class without an assigned successor is the last class that some professor teaches.

\section{Baseball Elimination}
\label{sec:11.6}

Every year millions of American baseball fans eagerly watch their favorite team, hoping they will win a spot in the playoffs, and ultimately the World Series.
Sadly, most teams are “mathematically eliminated” days or even weeks before the regular season ends.
Often, it is easy to spot when a team is eliminated—they can’t win enough games to catch up to the current leader in their division.
But sometimes the situation is more subtle.
For example, here are the actual standings from the American League East on August 30, 1996.
\[
\mbox{\sffamily\fontsize{9}{15}\selectfont\renewcommand{\arraystretch}{1.05}$
\begin{array}{c@{\hspace{1.4em}}c@{\hspace{1.1em}}c@{\hspace{0.8em}}|@{\hspace{0.8em}}c@{\hspace{0.9em}}c@{\hspace{0.9em}}c@{\hspace{0.9em}}c@{\hspace{0.9em}}c}
\hline
\mbox{\bfseries Team} & \mbox{{\bfseries Won}–{\bfseries Lost}} & \mbox{\bfseries Left} & \mbox{NYY} & \mbox{BAL} & \mbox{BOS} & \mbox{TOR} & \mbox{DET}\\
\hline
\mbox{New York Yankees} & \mbox{75–59} & \mbox{28} & & \mbox{3} & \mbox{8} & \mbox{7} & \mbox{3}\\
\mbox{Baltimore Orioles} & \mbox{71–63} & \mbox{28} & \mbox{3} & & \mbox{2} & \mbox{7} & \mbox{4}\\
\mbox{Boston Red Sox} & \mbox{69–66} & \mbox{27} & \mbox{8} & \mbox{2} & & \mbox{0} & \mbox{0}\\
\mbox{Toronto Blue Jays} & \mbox{63–72} & \mbox{27} & \mbox{7} & \mbox{7} & \mbox{0} & & \mbox{0}\\
\mbox{Detroit Tigers} & \mbox{49–86} & \mbox{27} & \mbox{3} & \mbox{4} & \mbox{0} & \mbox{0} & \\
\hline
\end{array}$}
\]

Detroit is clearly behind, but some die-hard Tigers fans may hold out hope that their team can still win.
After all, if Detroit wins all 27 of their remaining games, they will end the season with 76 wins, more than any other team has now.
So as long as every other team loses every game. . . but that’s not possible, because some of those other teams still have to play each other.
Here is one complete argument:\footnote{Both the example and this argument are taken from Eli Olinick’s web site https://s2.smu.edu/\textasciitilde olinick/riot/detroit.html, which is based on Olinick’s joint research with Ilan Adler, Alan Erera, and Dorit Hochbaum.}

\begin{quote}
\sffamily\fontsize{8.1}{11}\selectfont\setlength{\parindent}{1.6em}
By winning all of their remaining games, Detroit can finish the season with a record of 76 and 86. If the Yankees win just 2 more games, then they will finish the season with a 77 and 85 record which would put them ahead of Detroit. So, let’s suppose the Tigers go undefeated for the rest of the season and the Yankees fail to win another game.

The problem with this scenario is that New York still has 8 games left with Boston. If the Red Sox win all of these games, they will end the season with at least 77 wins putting them ahead of the Tigers. Thus, the only way for Detroit to even have a chance of finishing in first place, is for New York to win exactly one of the 8 games with Boston and lose all their other games. Meanwhile, the Sox must lose all the games they play against teams other than New York. This puts them in a 3-way tie for first place....

Now let’s look at what happens to the Orioles and Blue Jays in our scenario. Baltimore has 2 games left with with Boston and 3 with New York. So, if everything happens as described above, the Orioles will finish with at least 76 wins. So, Detroit can catch Baltimore only if the Orioles lose all their games to teams other than New York and Boston. In particular, this means that Baltimore must lose all 7 of its remaining games with Toronto. The Blue Jays also have 7 games left with the Yankees and we have already seen that for Detroit to finish in first place, Toronto must will all of these games. But if that happens, the Blue Jays will win at least 14 more games giving them at final record of 77 and 85 or better which means they will finish ahead of the Tigers. So, no matter what happens from this point in the season on, Detroit can not finish in first place in the American League East.
\end{quote}
\noindent There has to be a better way to figure this out!

Here is a more abstract formulation of the problem.
Our input consists of two arrays $W[1\,..\,n]$ and $G[1\,..\,n,1\,..\,n]$, where $W[i]$ is the number of games team $i$ has already won, and $G[i,j]$ is the number of upcoming games between teams $i$ and $j$.
We want to determine whether team $n$ can end the season with the most wins (possibly tied with other teams).\footnote{We are implicitly assuming that no game ends in a tie and that every game is actually played. Both assumptions are consistent with Major League Baseball rules, at least for games that affect postseason standing, barring wars, natural disasters, or swarms of bees.}

In the mid-1960s, Benjamin Schwartz observed that this question can be modeled as a maximum flow problem; about 20 years later, Dan Gusfield, Charles Martel, and David Fernández-Baca simplified Schwartz’s flow formulation to a pair selection problem.
Specifically, we want to know whether it is possible to \textbf{\textit{select}} \textit{a winner for each game}, so that team $n$ comes in first place.
Let $R[i] = \sum_j G[i,j]$ denote the number of remaining games for team $i$.
We will assume that team $n$ wins all $R[n]$ of its remaining games.
Then team $n$ can come in first place if and only if every other team $i$ wins at most $W[n] + R[n] - W[i]$ of its $R[i]$ remaining games.

Since we want to \textbf{\textit{select}} a winning team for each game, we start by building a bipartite graph, whose nodes represent the games and the teams.
We have $\binom{n}{2}$ \textit{game} nodes $g_{i,j}$, one for each pair $1 \le i < j < n$, and $n - 1$ \textit{team} nodes $t_i$, one for each $1 \le i < n$.
For each pair $i,j$, we add edges $\arc{g_{i,j}}{t_i}$ and $\arc{g_{i,j}}{t_j}$ with \textit{infinite} capacity.
We add a source vertex $s$ and edges $\arc{s}{g_{i,j}}$ with capacity $G[i,j]$ for each pair $i,j$.
Finally, we add a target node $t$ and edges $\arc{t_i}{t}$ with capacity $W[n] - W[i] + R[n]$ for each team $i$.

\figref{fig:11.7} shows the graph derived from the 1996 American League East standings, where “team $n$” is the Detroit Tigers.
All unlabeled edges have infinite capacity.

\begin{figure}[ht]
\centering
\includegraphics{p383-Im164}
\caption{Cubs win! Cubs win!}
\label{fig:11.7}
\end{figure}

\noindent\textbf{Theorem.} \textsl{Team $n$ can end the season in first place if and only if there is a feasible flow in this graph that saturates every edge leaving $s$.}

\noindent\textbf{Proof:} Suppose it is possible for team $n$ to end the season in first place.
Then every team $i < n$ wins at most $W[n] + R[n] - W[i]$ of the remaining games.
For each game between team $i$ and team $j$ that team $i$ wins, add one unit of flow along the path $\arc{\arc{\arc{s}{g_{i,j}}}{t_i}}{t}$.
Because there are exactly $G[i,j]$ games between teams $i$ and $j$, every edge leaving $s$ is saturated.
Because each team $i$ wins at most $W[n] + R[n] - W[i]$ games, the resulting flow is feasible.

Conversely, let $f$ be a feasible flow that saturates every edge out of $s$.
Suppose team $i$ wins exactly $f(\arc{g_{i,j}}{t_i})$ games against team $j$, for all $i$ and $j$.
Then teams $i$ and $j$ play $f(\arc{g_{i,j}}{t_i}) + f(\arc{g_{i,j}}{t_j}) = f(\arc{s}{g_{i,j}}) = G[i,j]$ games, so every upcoming game is played.
Moreover, each team $i$ wins a total of $\sum_j f(\arc{g_{i,j}}{t_i}) = f(\arc{t_i}{t}) \le W[n]+R[n]-W[i]$ upcoming games, and therefore at most $W[n] + R[n]$ games overall.
Thus, if team $n$ win all their upcoming games, they end the season in first place.\qed

In summary, to decide whether our favorite team can win, we construct the flow network, compute a maximum flow, and report whether than maximum flow saturates every edge out of $s$.
For example, in the graph in \figref{fig:11.7}, the total capacity of the edges leaving $s$ is 27 (because there are 27 remaining games).
On the other hand, the total capacity of the edges entering $t$ is only 26, which implies that the maximum flow value is at most 26.
We conclude that Detroit is mathematically eliminated.\footnote{We got (un)lucky here; it is possible for a team to be eliminated even if the total capacity of all edges into $t$ is no smaller than the total capacity of edges out of $s$.}

The flow network has $O(n^2)$ vertices and $O(n^2)$ edges, and it can be constructed in $O(n^2)$ time.
Using Orlin’s algorithm, we can compute the maximum flow in $O(VE) = \symbfit{O(n^4)}$ \textbf{\textit{time}}.

This is not the fastest algorithm for the baseball elimination problem.
In 2001, Kevin Wayne proved that one can determine \textbf{\textit{all}} teams that are mathematically eliminated in only $\symbfit{O(n^3)}$ \textbf{\textit{time}}, essentially using a single maximum-flow computation.

\section{Project Selection}
\label{sec:11.7}

In our final example, suppose we are given a set of $n$ projects that we could possibly perform.
Some projects cannot be started until certain other projects are completed.
The projects and their dependencies are described by a directed acyclic graph $G$ whose vertices are the projects, where each edge $\arc{u}{v}$ indicates that project $u$ cannot be performed before project $v$.
(This is exactly the form of dependency graphs we considered in \chapref{sec:6.4}.)
Finally, each project $v$ has an associated \textit{profit} $\$(v)$ which will be given to us if the project is completed; some projects have negative profits, which we interpret as positive \textit{costs}.
We can choose to finish any subset $X$ of the projects that includes all its dependents; that is, for every project $x \in X$, every project that $x$ depends on is also in $X$.
Our goal is to find a valid subset of the projects whose total profit is as large as possible.
In particular, if all of the jobs have negative profit, the correct answer is to do nothing.

\begin{figure}[ht]
\centering
\includegraphics{p384-Im165}
\caption{A dependency graph for a set of eight projects. Diamonds indicate profitable projects; squares indicate costly projects. Each edge $\arc{u}{v}$ means $u$ depends on $v$.}
\label{fig:11.8}
\end{figure}

At a high level, our task to partition the projects into two subsets $S$ and $T$, the jobs we \textit{Select} and the jobs we \textit{Turn down}.
So intuitively, we’d like to model our problem as a minimum cut problem in a certain graph.
But in which graph?
How do we enforce prerequisites?
We want to \textit{maximize} profit, but we only know how to find \textit{minimum} cuts.
And how do we convert negative profits into positive capacities?

To transform our given constraint graph $G$ into a flow network $G'$, we add a source vertex $s$ and a target vertex $t$ to the dependency graph, with an edge $\arc{s}{v}$ for every profitable job $v$ (with $\$(v) > 0$), and an edge $\arc{u}{t}$ for every costly job $u$ (with $\$(u) < 0$).
Intuitively, we can think of $s$ as a new job (“Sleep!”) with profit/cost $0$ that we must perform last.
We assign capacities to the edges of $G'$ as follows:
\begin{itemize}
\item $c(\arc{s}{v}) = \$(v)$ for every profitable job $v$;
\item $c(\arc{u}{t}) = -\$(u)$ for every costly job $u$;
\item $c(\arc{u}{v}) = \infty$ for every dependency edge $\arc{u}{v}$.
\end{itemize}
\noindent All edge-capacities are positive, so this is a valid input to the maximum cut problem.

Now consider an arbitrary $(s,t)$-cut $(S,T)$ in $G'$.
For any edge $\arc{u}{v}$ in the original dependency graph, if $u \in S$ and $v \in T$, then $\lVert S,T\rVert = \infty$.
Thus, we can legally select the jobs in $S$ if and only if the capacity of the cut $(S,T)$ is finite.

\begin{figure}[ht]
\centering
\includegraphics{p385-Im166}
\caption{The flow network for the example dependency graph, along with its minimum cut. The cut has capacity $13$ and $P = 15$, so the total profit for the selected jobs is $2$.}
\label{fig:11.9}
\end{figure}

In fact, it turns out that cuts with smaller capacity correspond to job selections with higher profit.
Specifically, I claim that selecting the jobs in $S$ earns a total profit of $P - \lVert S,T\rVert$, where $P$ is the sum of all the positive profits:
\[
P = \sum_v \max\{0, \$(v)\} = \sum_{\$(v)>0} \$(v).
\]
We can prove this claim by straightforward definition-chasing, as follows.
For any subset $X$ of projects, we define three values.
(Here, as usual, we define $c(\arc{u}{v}) = 0$ when $\arc{u}{v}$ is not an edge.)
\begin{align*}
\Var{cost}(X) &:= \sum_{\substack{u \in X\\ \$(u)<0}} -\$(u) \;=\; \sum_{u \in X} c(\arc{u}{t})\\
\Var{yield}(X) &:= \sum_{\substack{v \in X\\ \$(v)>0}} \$(v) \;=\; \sum_{v \in X} c(\arc{s}{v})\\
\Var{profit}(X) &:= \sum_{v \in X} \$(v) \;=\; \Var{yield}(X) - \Var{cost}(X).
\end{align*}
By definition, $P = \Var{yield}(V) = \Var{yield}(S) + \Var{yield}(T)$.
Because the cut $(S,T)$ has finite capacity, only edges of the form $\arc{s}{v}$ and $\arc{u}{t}$ can cross the cut.
By construction, every edge $\arc{s}{v}$ points to a profitable job and each edge $\arc{u}{t}$ points from a costly job.
Thus, $\lVert S,T\rVert = \Var{cost}(S) + \Var{yield}(T)$.
We immediately conclude that $P - \lVert S,T\rVert = \Var{yield}(S) - \Var{cost}(S) = \Var{profit}(S)$, as claimed.

It follows immediately that we can \textit{maximize} our total profit by computing a \textit{minimum} cut in $G'$.
We can easily construct $G'$ from $G$ in $O(V + E)$ time, and we can compute the minimum $(s,t)$-cut in $G'$ in $O(VE)$ time using Orlin’s algorithm.
We conclude that the entire project-selection algorithm runs in $\symbfit{O(VE)}$ \textbf{\textit{time}}.

\begin{exercises}
\item
\label{ex:11.1}
Let $G = (V,E)$ be a directed graph where for each vertex $v$, the in-degree and out-degree of $v$ are equal.
Suppose $G$ contains $k$ edge-disjoint paths from some vertex $u$ to another vertex $v$.
Under these conditions, must $G$ also contain $k$ edge-disjoint paths from $v$ to $u$?
Give a proof or a counterexample with explanation.

\item
\label{ex:11.2}
Given an undirected graph $G = (V,E)$, with three vertices $u$, $v$, and $w$, describe and analyze an algorithm to determine whether there is a path from $u$ to $w$ that passes through $v$.
\Hint{If $G$ were a directed graph, this problem would be NP-hard!}

\item
\label{ex:11.3}
Consider a directed graph $G = (V,E)$ with several source vertices $s_1, s_2, \dots, s_\sigma$ and target vertices $t_1, t_1, \dots, t_\tau$, where no vertex is both a source and a target.
A \textit{multi-terminal flow} is a function $f : E \to \mathbb{R}_{\ge 0}$ that satisfies the flow conservation constraint at every vertex that is neither a source nor a target.
The value $|f|$ of a multi-terminal flow is the total excess flow out of \textit{all} the source vertices:
\[
|f| := \sum_{i=1}^{\sigma} \left( \sum_w f(\arc{s_i}{w}) - \sum_u f(\arc{u}{s_i}) \right)
\]
As usual, we are interested in finding flows with maximum value, subject to capacity constraints on the edges.
(In particular, we don’t care how much flow moves from any particular source to any particular target.)
\begin{enumerate}
\item Consider the following algorithm for computing multi-terminal flows.
The variables $f$ and $f'$ represent flow functions.
The subroutine \Proc{MaxFlow}$(G, s, t)$ solves the standard maximum flow problem with source $s$ and target $t$.
\begin{pseudocode}
\Header{\Proc{MaxMultiFlow}$(G, s[1\,..\,\sigma], t[1\,..\,\tau])$:}\\
$f \gets 0$ \Comment{Initialize the flow}\\
for $i \gets 1$ to $\sigma$\\
\Indent for $j \gets 1$ to $\tau$\\
\Indent\Indent $f' \gets \Proc{MaxFlow}(G_f, s[i], t[j])$\\
\Indent\Indent $f \gets f + f'$ \Comment{Update the flow}\\
return $f$
\end{pseudocode}
\noindent Prove that this algorithm correctly computes a maximum multi-terminal flow in $G$.
\item Describe a more efficient algorithm to compute a maximum multi-terminal flow in $G$.
\end{enumerate}

\item
\label{ex:11.4}
The Island of Sodor is home to a large number of towns and villages, connected by an extensive rail network.
Recently, several cases of a deadly contagious disease (either swine flu or zombies; reports are unclear) have been reported in the village of Skarloey.
The controller of the Sodor railway plans to close down certain railway stations to prevent the disease from spreading to Tidmouth, his home town.
No trains can pass through a closed station.
To minimize expense (and public notice), he wants to close down as few stations as possible.
However, he cannot close the Skarloey station, because that would expose him to the disease, and he cannot close the Tidmouth station, because then he couldn’t visit his favorite pub.

Describe and analyze an algorithm to find the minimum number of stations that must be closed to block all rail travel from Skarloey to Tidmouth.
The Sodor rail network is represented by an undirected graph, with a vertex for each station and an edge for each rail connection between two stations.
Two special vertices $s$ and $t$ represent the stations in Skarloey and Tidmouth.

For example, given the following input graph, your algorithm should return the integer $2$.
\begin{center}
\includegraphics{p387-Im167}
\end{center}

\item
\label{ex:11.5}
An $n \times n$ grid is an undirected graph with $n^2$ vertices organized into $n$ rows and $n$ columns.
We denote the vertex in the $i$th row and the $j$th column by $(i,j)$.
Every vertex $(i,j)$ has exactly four neighbors $(i - 1, j)$, $(i + 1, j)$, $(i, j - 1)$, and $(i, j + 1)$, except the \textit{boundary} vertices, for which $i = 1$, $i = n$, $j = 1$, or $j = n$.

Let $(x_1, y_1), (x_2, y_2), \dots, (x_m, y_m)$ be distinct vertices, called \textit{terminals}, in the $n \times n$ grid.
The \textbf{\textit{escape problem}} is to determine whether there are $m$ vertex-disjoint paths in the grid that connect the terminals to any $m$ distinct boundary vertices.
\begin{enumerate}
\item Describe and analyze an efficient algorithm to solve the escape problem.
The running time of your algorithm should be a small polynomial function of $n$.
\item Now suppose the input to the escape problem consists of a single integer $n$ and the list of $m$ terminal vertices.
If $m$ is very small, the previous running time is actually \textit{exponential} in the input size!
Describe and analyze an algorithm to solve the escape problem in time \textit{polynomial in} $m$.
\item \difficulty{h} Modify the previous algorithm to output an explicit description of the escape paths (if they exist), still in time polynomial in $m$.
\end{enumerate}

\begin{figure}[ht]
\centering
\includegraphics{p388-Im168}
\caption{A positive instance of the escape problem, and its solution.}
\label{fig:11.10}
\end{figure}

\item
\label{ex:11.6}
The SPU Commuter Silence Department is installing a mini-golf course in the basement of the See-Bull Center!
The playing field is a closed polygon bounded by $m$ horizontal and vertical line segments, meeting at right angles.
The course has $n$ \textit{starting points} and $n$ \textit{holes}, in one-to-one correspondence.
It is always possible hit the ball along a straight line directly from each starting point to the corresponding hole, without touching the boundary of the playing field.
(Players are not allowed to bounce golf balls off the walls; too much glass.)
The $n$ starting points and $n$ holes are all at distinct locations.

Sadly, the architect’s computer crashed just as construction was about to begin.
Thanks to the herculean efforts of their sysadmins, they were able to recover the \textit{locations} of the starting points and the holes, but all information about which starting points correspond to which holes was lost!

Describe and analyze an algorithm to compute a one-to-one correspondence between the starting points and the holes that meets the straight-line requirement, or to report that no such correspondence exists.
The input consists of the $x$- and $y$-coordinates of the $m$ corners of the playing field, the $n$ starting points, and the $n$ holes.
Assume you can determine in constant time whether two line segments intersect, given the $x$- and $y$-coordinates of their endpoints.

\begin{figure}[ht]
\centering
\includegraphics{p389-Im169}
\caption{A mini-golf course with five starting points ($\star$) and holes ($\smwhtcircle$), and a legal correspondence between them.}
\label{fig:11.11}
\end{figure}

\item
\label{ex:11.7}
A \textbf{\textit{cycle cover}} of a given directed graph $G = (V,E)$ is a set of vertex-disjoint cycles that cover every vertex in $G$.
Describe and analyze an efficient algorithm to find a cycle cover for a given graph, or correctly report that no cycle cover exists.
\Hint{Use bipartite matching!}

\item
\label{ex:11.8}
Suppose you are given an $n \times n$ checkerboard with some of the squares deleted.
You have a large set of dominos, just the right size to cover two squares of the checkerboard.
Describe and analyze an algorithm to determine whether one tile the board with dominos—each domino must cover exactly two undeleted squares, and each undeleted square must be covered by exactly one domino.

Your input is a boolean array $\Var{Deleted}[1\,..\,n,1\,..\,n]$, where $\Var{Deleted}[i,j] = \Const{True}$ if and only if the square in row $i$ and column $j$ has been deleted.
Your output is a single boolean; you do \textbf{\textit{not}} have to compute the actual placement of dominos.
For example, for the board shown in \figref{fig:11.12}, your algorithm should return \Const{True}.

\begin{figure}[ht]
\centering
\includegraphics{p389-Im170}
\caption{Covering a partial checkerboard with dominos.}
\label{fig:11.12}
\end{figure}

\item
\label{ex:11.9}
Suppose we are given an $n \times n$ square grid, some of whose squares are colored black and the rest white.
Describe and analyze an algorithm to determine whether tokens can be placed on the grid so that
\begin{itemize}
\item every token is on a white square;
\item every row of the grid contains exactly one token; and
\item every column of the grid contains exactly one token.
\end{itemize}

Your input is a two dimensional array $\Var{IsWhite}[1\,..\,n,1\,..\,n]$ of booleans, indicating which squares are white.
Your output is a single boolean.
For example, given the grid in \figref{fig:11.13} as input, your algorithm should return \Const{True}.

\begin{figure}[ht]
\centering
\includegraphics{p390-Im171}
\caption{Marking every row and column in a grid.}
\label{fig:11.13}
\end{figure}

\item
\label{ex:11.10}
Suppose we are given a set of boxes, each specified by their height, width, and depth in centimeters.
All three side lengths of every box lie strictly between 10cm and 20cm.
As you should expect, one box can be placed inside another if the first box can be rotated so that its height, width, and depth are respectively smaller than the height, width, and depth of the second box.
Boxes can be nested recursively.
Call a box is \textit{visible} if it is not inside another box.

Describe and analyze an algorithm to nest the boxes so that the number of visible boxes is as small as possible.

\item
\label{ex:11.11}
Suppose we are given an $n \times n$ grid, some of whose cells are marked; the grid is represented by an array $M[1\,..\,n,1\,..\,n]$ of booleans, where $M[i,j] = \Const{True}$ if and only if cell $(i,j)$ is marked.
A \textit{monotone} path through the grid starts at the top-left cell, moves only right or down at each step, and ends at the bottom-right cell.
Our goal is to cover the marked cells with as few monotone paths as possible.
\begin{enumerate}
\item Describe an algorithm to find a monotone path that covers the largest number of marked cells.
\item There is a natural greedy heuristic to find a small cover by monotone paths: If there are any marked cells, find a monotone path $\pi$ that covers the largest number of marked cells, unmark any cells covered by $\pi$ those marked cells, and recurse.
Show that this algorithm does \textit{not} always compute an optimal solution.
\item Describe and analyze an efficient algorithm to compute the smallest set of monotone paths that covers every marked cell.
\end{enumerate}

\begin{figure}[ht]
\centering
\includegraphics{p390-Im172}
\caption{Greedily covering the marked cells in a grid with four monotone paths.}
\label{fig:11.14}
\end{figure}

\item
\label{ex:11.12}
The Faculty Senate at Sham-Poobanana University has decided to convene a committee to determine whether Uncle Gabby, Professor Bobo Cornelius, or Mofo the Psychic Gorilla should replace the recently disgraced Baron Factotum as the new official \Strike{mascot} \textit{symbol} of SPU’s athletic teams (The Fighting Pooh-bahs).
Exactly one faculty member must be chosen from each academic department to serve on this committee.
Some faculty members have appointments in multiple departments, but each committee member can represent only one department.
For example, if Prof. Blagojevich is affiliated with both the Department of Corruption and the Department of Stupidity, and he is chosen as the Stupidity representative, then someone else must represent Corruption.
Finally, University policy requires that every faculty committee must contain exactly the same number of assistant professors, associate professors, and full professors.
Fortunately, the number of departments is a multiple of 3.

Describe and analyze an algorithm to choose a subset of the SPU faculty to staff The Post-Factotum Simian \Strike{Mascot} \textit{Symbol} Committee, or correctly report that no valid committee is possible.
Your input is a bipartite graph indicating which professors belong to which departments; each professor vertex is labeled with that professor’s rank (assistant, associate, or full).

\item
\label{ex:11.13}
The Department of Commuter Silence at Sham-Poobanana University has a flexible curriculum with a complex set of graduation requirements.
The department offers $n$ different courses, and there are $m$ different requirements.
Each requirement specifies a subset of the $n$ courses and the number of courses that must be taken from that subset.
The subsets for different requirements may overlap, but each course can be used to satisfy \textit{at most one} requirement.

For example, suppose there are $n = 5$ courses $A, B, C, D, E$ and $m = 2$ graduation requirements:
\begin{itemize}
\item You must take at least $2$ courses from the subset $\{A, B, C\}$.
\item You must take at least $2$ courses from the subset $\{C, D, E\}$.
\end{itemize}
Then a student who has only taken courses $B, C, D$ cannot graduate, but a student who has taken either $A, B, C, D$ or $B, C, D, E$ can graduate.

Describe and analyze an algorithm to determine whether a given student can graduate.
The input to your algorithm is the list of $m$ requirements (each specifying a subset of the $n$ courses and the number of courses that must be taken from that subset) and the list of courses the student has taken.

\item
\label{ex:11.14}
You’re organizing the First Annual SPU Commuter Silence 72-Hour Dance Exchange, to be held all day Friday, Saturday, and Sunday.
Several 30-minute sets of music will be played during the event, and a large number of DJs have applied to perform.
You need to hire DJs according to the following constraints.
\begin{itemize}
\item Exactly $k$ sets of music must be played each day, and thus $3k$ sets altogether.
\item Each set must be played by a single DJ in a consistent music genre (ambient, bubblegum, dubstep, horrorcore, K-pop, Kwaito, mariachi, straight-ahead jazz, trip-hop, Nashville country, parapara, ska, \dots).
\item Each genre must be played at most once per day.
\item Each candidate DJ has given you a list of genres they are willing to play.
\item Each DJ can play at most three sets during the entire event.
\end{itemize}
Suppose there are $n$ candidate DJs and $g$ different musical genres available.
Describe and analyze an efficient algorithm that either assigns a DJ and a genre to each of the $3k$ sets, or correctly reports that no such assignment is possible.

\item
\label{ex:11.15}
Suppose you are running a web site that is visited by the same set of people every day.
Each visitor claims membership in one or more \textit{demographic groups}; for example, a visitor might describe himself as male, 40--50 years old, a father, a resident of Illinois, an academic, a blogger, and a fan of Gilbert and Sullivan.\footnote{I am a very good theoretical computer scientist, specifically, a geometric algorithm specialist.}
Your site is supported by advertisers.
Each advertiser has told you which demographic groups should see its ads and how many of its ads you must show each day.
Altogether, there are $n$ visitors, $k$ demographic groups, and $m$ advertisers.

Describe an efficient algorithm to determine, given all the data described in the previous paragraph, whether you can show each visitor exactly \textit{one} ad per day, so that every advertiser has its desired number of ads displayed, and every ad is seen by someone in an appropriate demographic group.

\item
\label{ex:11.16}
Suppose we are given an array $A[1\,..\,m][1\,..\,n]$ of non-negative real numbers.
We want to \textit{round} $A$ to an integer matrix, by replacing each entry $x$ in $A$ with either $\lfloor x \rfloor$ or $\lceil x \rceil$, without changing the sum of entries in any row or column of $A$.
For example:
\[
\begin{bmatrix}
1.2 & 3.4 & 2.4\\
3.9 & 4.0 & 2.1\\
7.9 & 1.6 & 0.5
\end{bmatrix}
\longmapsto
\begin{bmatrix}
1 & 4 & 2\\
4 & 4 & 2\\
8 & 1 & 1
\end{bmatrix}
\]
\begin{enumerate}
\item Describe and analyze an efficient algorithm that either rounds $A$ in this fashion, or reports correctly that no such rounding is possible.
\item Prove that a legal rounding is possible \textit{if and only if} the sum of entries in each row is an integer, and the sum of entries in each column is an integer.
In other words, prove that either your algorithm from part (a) returns a legal rounding, or a legal rounding is \textit{obviously} impossible.
\item \difficulty{h} Suppose we are guaranteed that none of the entries in the input matrix $A$ is an integer.
Describe and analyze an even faster algorithm that either rounds $A$ or reports correctly that no such rounding is possible.
For full credit, your algorithm must run in $O(mn)$ time.
\Hint{\textbf{Don’t} use flows.}
\end{enumerate}

\item
\label{ex:11.17}
\textbf{\textit{Ad-hoc}} \textbf{\textit{networks}} are made up of low-powered wireless devices.
In principle,\footnote{but not so much in practice} these networks can be used on battlefields, in regions that have recently suffered from natural disasters, and in other hard-to-reach areas.
The idea is that a large collection of cheap, simple devices could be distributed through the area of interest (for example, by dropping them from an airplane); the devices would then automatically configure themselves into a functioning wireless network.

These devices can communicate only within a limited range.
We assume all the devices are identical; there is a distance $D$ such that two devices can communicate if and only if the distance between them is at most $D$.

We would like our ad-hoc network to be reliable, but because the devices are cheap and low-powered, they frequently fail.
If a device detects that it is likely to fail, it should transmit its information to some other \textit{backup} device within its communication range.
We require each device $x$ to have $k$ potential backup devices, all within distance $D$ of $x$; we call these $k$ devices the \textbf{\textit{backup set}} of $x$.
Also, we do not want any device to be in the backup set of too many other devices; otherwise, a single failure might affect a large fraction of the network.

So suppose we are given the communication radius $D$, parameters $b$ and $k$, and an array $d[1\,..\,n,1\,..\,n]$ of distances, where $d[i,j]$ is the distance between device $i$ and device $j$.
Describe an algorithm that either computes a backup set of size $k$ for each of the $n$ devices, such that no device appears in more than $b$ backup sets, or reports (correctly) that no good collection of backup sets exists.

\item
\label{ex:11.18}
Faced with the threat of brutally severe budget cuts, Potemkin University has decided to hire actors to sit in classes as “students”, to ensure that every class they offer is completely full.
Because actors are expensive, the university wants to hire as few of them as possible.

Building on their previous leadership experience at the now-defunct Sham-Poobanana University, the administrators at Potemkin have given you a directed acyclic graph $G = (V,E)$, whose vertices represent classes, and where each edge $\arc{i}{j}$ indicates that the same “student” can attend class $i$ and then later attend class $j$.
In addition, you are also given an array $\Var{cap}[1\,..\,V]$ listing the maximum number of “students” who can take each class.
Describe an analyze an algorithm to compute the minimum number of “students” that would allow every class to be filled to capacity.

\item
\label{ex:11.19}
Quentin, Alice, and the other Brakebills Physical Kids are planning an excursion through the Neitherlands to Fillory.
The Neitherlands is a vast, deserted city composed of several plazas, each containing a single fountain that can magically transport people to a different world.
Adjacent plazas are connected by gates, which have been cursed by the Beast.
The gates between plazas are open only for five minutes every hour, all simultaneously—from 12:00 to 12:05, then from 1:00 to 1:05, and so on—and are otherwise locked.
During those five minutes, if more than one person passes through any single gate, the Beast will detect their presence.\footnote{This is very bad.}
Moreover, anyone attempting to open a locked gate, or attempting to pass through more than one gate within the same five-minute period will turn into a niffin.\footnote{This is very very bad.}
However, any number of people can safely pass through \textit{different} gates at the same time and/or pass through the same gate at \textit{different} times.

You are given a map of the Neitherlands, which is a graph $G$ with a vertex for each fountain and an edge for each gate, with the fountains to Earth and Fillory clearly marked.
\begin{enumerate}
\item Suppose you are also given a positive integer $h$.
Describe and analyze an algorithm to compute the maximum number of people that can walk from the Earth fountain to the Fillory fountain in at most $h$ hours—that is, after the gates have opened at most $h$ times—without anyone alerting the Beast or turning into a niffin.
The running time of your algorithm should depend on $h$.
\Hint{Build a different graph.}
\item \difficulty{ch} Describe an analyze an algorithm for part (a) whose running time is polynomial in $V$ and $E$, with \textit{no} dependence on $h$.
\item On the other hand, suppose you are also given an integer $k$.
Describe and analyze an algorithm to compute the minimum number of hours that allow $k$ people to walk from the Earth fountain to the Fillory fountain, without anyone alerting the Beast or turning into a niffin.
\Hint{Use part (a).}
\end{enumerate}

\item \difficulty{h}
\label{ex:11.20}
Let $G = (L \sqcup R, E)$ be a bipartite graph, whose left vertices $L$ are indexed $\ell_1, \ell_2, \dots, \ell_n$ and whose right vertices are indexed $r_1, r_2, \dots, r_n$.
A matching $M$ in $G$ is \textbf{\textit{non-crossing}} if, for every pair of edges $\ell_i r_j$ and $\ell_{i'} r_{j'}$ in $M$, we have $i < i'$ if and only if $j < j'$.
\begin{enumerate}
\item Describe and analyze an algorithm to find the largest non-crossing matching in $G$.
\Hint{This is not really a flow problem.}
\item Describe and analyze an algorithm to find the smallest number of non-crossing matchings $M_1, M_2, \dots, M_k$ such that each edge in $G$ is in exactly one matching $M_i$.
\Hint{This is really a flow problem.}
\end{enumerate}

\item \difficulty{h}
\label{ex:11.21}
Let $G = (L \sqcup R, E)$ be a bipartite graph, whose left vertices $L$ are indexed $\ell_1, \ell_2, \dots, \ell_n$ in some arbitrary order.
\begin{enumerate}
\item A matching $M$ in $G$ is \textbf{\textit{dense}} if there are no consecutive unmatched vertices in $L$; that is, for each index $i$, at least one of the vertices $\ell_i$ and $\ell_{i+1}$ is incident to an edge in $M$.
Describe an algorithm to determine whether $G$ has a dense matching.
\item A matching $M$ in $G$ is \textbf{\textit{sparse}} if there are no consecutive \textit{matched} vertices in $L$; that is, for each index $i$, at least one of the vertices $\ell_i$ and $\ell_{i+1}$ is \textit{not} incident to an edge in $M$.
(In particular, the empty matching is sparse.)
Describe an algorithm to find the largest sparse matching in $G$.
\item A matching $M$ in $G$ is \textbf{\textit{palindromic}} if, for every index $i$, either $\ell_i$ and $\ell_{n-i+1}$ are both incident to edges in $M$, or neither $\ell_i$ nor $\ell_{n-i+1}$ is incident to an edge in $M$.
(In particular, the empty matching is palindromic.)
Describe an algorithm to find the largest palindromic matching in $G$.
\end{enumerate}
None of these problems restrict which vertices in $R$ are matched or unmatched.

\item \difficulty{h}
\label{ex:11.22}
A \textit{rooted tree} is a directed acyclic graph, in which every vertex has exactly one incoming edge, except for the \textit{root}, which has no incoming edges.
Equivalently, a rooted tree consists of a root vertex, which has edges pointing to the roots of zero or more smaller rooted trees.
Describe an efficient algorithm to compute, given two rooted trees $A$ and $B$, the largest rooted tree that is isomorphic to both a subgraph of $A$ and a subgraph of $B$.
More briefly, describe an algorithm to find the largest common subtree of two rooted trees.

\Hint{This would be a relatively straightforward dynamic programming problem if either every node had $O(1)$ children or the children of each node were ordered from left to right. But for unordered trees with large degree, you need another technique to combine recursive subproblems efficiently.}
\end{exercises}
